\documentclass[conference]{IEEEtran}

\usepackage[english]{babel}
\usepackage{amsmath}
\usepackage{graphicx}
\usepackage{url}
\usepackage{microtype}

\title{Janus: Evidence-Based Governance for AI-Assisted Development Systems}

\author{%
  \IEEEauthorblockN{Martín Nicolás Sánchez Morales}\\
  \IEEEauthorblockA{Janus Governance Project\\
  Argentina}%
}

\begin{document}

\maketitle

\begin{abstract}
AI-assisted development introduces structural challenges in traceability, authority,
and reproducibility. This paper presents Janus, a protocol-driven governance layer
that records and signals deviations from evidence-based workflow protocols across
development environments.

Rather than measuring system performance, this work evaluates governance through
conformance between protocol-defined expectations and runtime-observed behavior.

We introduce an evidence-based governance benchmark and present preliminary
validation results across multiple runtime artifacts.
\end{abstract}

\begin{IEEEkeywords}
Governance, Observability, AI-assisted development, Traceability, Runtime systems
\end{IEEEkeywords}

\section{Introduction}

AI-assisted systems introduce opacity into decision-making processes, making it
difficult to trace actions, verify authority, and reconstruct execution history.

Janus introduces a governance layer that records and signals deviations from
protocol-defined workflows, providing:
\begin{itemize}
    \item Explicit evidence (E\textsuperscript{+} and E\textsuperscript{-})
    \item Human authority boundaries
    \item Deterministic reconstruction of governance state
\end{itemize}

Rather than modifying underlying systems, Janus wraps workflows with a protocol
stack that monitors traceability and signals accountability deviations at the operational level.

\section{Architecture}

Janus operates as a protocol-driven governance layer composed of four distinct
components:

\begin{itemize}
    \item \textbf{Core Layer:} defines events, evidence types, and authority boundaries
    \item \textbf{Protocol Layer:} defines operational governance rules
    \item \textbf{Runtime Layer:} captures telemetry and execution events
    \item \textbf{Workflow Layer:} mediates interaction between agents and human authority
\end{itemize}

\section{Governance Invariants}

Janus defines four operational invariants that govern system behavior:

\medskip
\noindent\textbf{Invariant 1 (Evidence Completeness):}
Every governance-relevant action must produce either E\textsuperscript{+}
(recorded occurrence) or E\textsuperscript{-} (recorded omission).

\medskip
\noindent\textbf{Invariant 2 (Omission Detectability):}
Every omission must be detectable under a deterministic, predefined scope.

\medskip
\noindent\textbf{Invariant 3 (Human Authority Closure):}
No governance state is considered complete without an explicit human decision
or a recorded omission event.

\medskip
\noindent\textbf{Invariant 4 (Deterministic Reconstruction):}
Any given sequence of governance events must produce the same governance state
upon reconstruction.

\section{Runtime Observability}

Janus emits structured events at runtime, each containing:
\begin{itemize}
    \item Event type
    \item Timestamp
    \item Payload
    \item Error traces
\end{itemize}

This structured telemetry enables immediate failure detection and end-to-end
traceability. The observed GVL values (95--120\,ms) indicate near-immediate
detection under interactive system conditions.

\section{Evidence-Based Governance Benchmark}

\subsection{Metric Definitions}

\noindent\textbf{Event Capture Rate (ECR):}
\begin{equation}
  ECR = \frac{\text{Governed Events}}{\text{Total System Operations}}
\end{equation}
ECR measures the proportion of system operations covered by governance
instrumentation. A value of 1.0 indicates full governance coverage of all
observed operations.

\medskip
\noindent\textbf{Governance Verdict Latency (GVL):}
\begin{equation}
\begin{split}
  GVL = \; & t(\mathit{HUMAN\_VERDICT}) \\
            & - \; t(\mathit{OMISSION\_DETECTED})
\end{split}
\end{equation}
GVL is measured in milliseconds and reflects the elapsed time between
omission detection and the recording of a human authority verdict.

\medskip
\noindent\textbf{Protocol Violation Detectability Rate (PVDR):}
\begin{equation}
  PVDR = \frac{\text{Detected Violations}}{\text{Injected Violations}}
\end{equation}

\subsection{Evaluation Principle}

\begin{equation}
  \mathit{Governance\_Conformance} =
    \frac{\text{Observed Events}}{\text{Expected Protocol Events}}
\end{equation}

Expected protocol events are derived from predefined protocol rules that specify
required event types under given execution conditions.

\section{Experimental Validation}

\subsection{Fault Injection Method}

Protocol violations were introduced through controlled test cases, including:
\begin{itemize}
    \item Missing required decision events
    \item Invalid state transitions
    \item Incomplete evidence generation
\end{itemize}
A total of 50 injected violations per artifact were used.
Human verdict events were simulated through controlled test inputs, allowing
consistent timestamp capture for GVL measurement under experimental conditions.

\subsection{Results}

Validation was performed across multiple runtime artifacts under controlled
development conditions.

\begin{table}[h]
  \centering
  \caption{Preliminary Benchmark Results}
  \label{tab:results}
  \begin{tabular}{|l|c|c|c|}
    \hline
    \textbf{Artifact} & \textbf{ECR} & \textbf{GVL (ms)} & \textbf{PVDR} \\
    \hline
    Interactive System & 1.0 & 120 & 0.92 \\
    Logic Component    & 1.0 &  95 & 0.95 \\
    Rendering Module   & 1.0 & 110 & 0.90 \\
    \hline
  \end{tabular}
\end{table}

These results represent initial validation under controlled development conditions.

\section{Related Work}

Janus relates to several established paradigms and frameworks:

\begin{itemize}
    \item Event Sourcing~\cite{fowler}
    \item Command Query Responsibility Segregation (CQRS)~\cite{young}
    \item Distributed event ordering~\cite{lamport}
    \item Observability frameworks~\cite{otel}
    \item AI risk management~\cite{nist}
    \item Requirements traceability~\cite{gotel}
\end{itemize}

Janus differs from these approaches by integrating governance, evidence, and
human authority as first-class runtime constructs, rather than as post-hoc
auditing mechanisms.

\section{Limitations}

Current limitations of the Janus model include:

\begin{itemize}
    \item Lack of automated enforcement of authority boundaries
    \item Partial observability in external or third-party systems
    \item Human authority scalability challenges at high decision volumes
\end{itemize}

The system currently focuses on detection and signaling of governance deviations,
rather than enforcing them automatically.

\subsection{Threats to Validity}
\begin{itemize}
    \item Limited number of runtime artifacts evaluated
    \item Single-environment validation scope
    \item Lack of independent replication by external teams
    \item Controlled fault injection scenarios may not reflect production conditions
\end{itemize}

\section{Conclusion}

Janus demonstrates that governance can be evaluated through measurable conformance
between expected and observed behavior in runtime systems. By treating evidence
capture, omission detection, and human authority as first-class operational
constructs, Janus provides a reproducible foundation for governance in AI-assisted
development environments.

\begin{thebibliography}{1}

\bibitem{fowler}
M. Fowler, ``Event Sourcing,'' 2005. [Online].
Available: \url{https://martinfowler.com/eaaDev/EventSourcing.html}

\bibitem{young}
G. Young, ``CQRS Documents,'' 2010. [Online].
Available: \url{https://cqrs.files.wordpress.com/2010/11/cqrs_documents.pdf}

\bibitem{lamport}
L. Lamport, ``Time, Clocks, and the Ordering of Events in a Distributed System,''
\emph{Commun. ACM}, vol.~21, no.~7, pp.~558--565, Jul. 1978.

\bibitem{otel}
OpenTelemetry, ``OpenTelemetry Specification,''
Cloud Native Computing Foundation (CNCF), 2023. [Online].
Available: \url{https://opentelemetry.io/docs/specs/otel/}

\bibitem{nist}
NIST, ``AI Risk Management Framework (AI RMF 1.0),'' 2023. [Online].
Available: \url{https://www.nist.gov/itl/ai-risk-management-framework}

\bibitem{gotel}
O. Gotel and A. Finkelstein, ``An Analysis of the Requirements Traceability
Problem,'' \emph{Proc. IEEE Int. Conf. Requirements Engineering}, 1994.

\end{thebibliography}

\end{document}
